%%%%%%%%%%%%%%%%%%%%%%%%%%%%%%%%%%%%%%%%%%%%%%%%%%%%%%%%%%%%%%%%%%%%%%%%%%%%%%%%
%                                                                              %
% BlackArch Linux Guide                                                        %
%                                                                              %
%%%%%%%%%%%%%%%%%%%%%%%%%%%%%%%%%%%%%%%%%%%%%%%%%%%%%%%%%%%%%%%%%%%%%%%%%%%%%%%%


%%% BEGIN %%%
\documentclass[a4paper, oneside, 11pt]{book}


%%% INCLUDES %%%
\renewcommand{\familydefault}{\sfdefault}
\usepackage{verbatim}
\usepackage[utf8]{inputenc}
\usepackage{geometry}
\usepackage{array}
\usepackage{graphicx}
\usepackage{supertabular}
\usepackage{verbatim}
\usepackage{pstricks}
\usepackage{fancyhdr}
\usepackage{fancyvrb}
\usepackage{tocloft}
\usepackage{color}
\usepackage{enumerate}
\usepackage[french]{babel}
\usepackage[
pdftitle={BlackArch Linux, The BlackArch Linux Guide},
pdfsubject={BlackArch Linux, The BlackArch Linux Guide},
pdfauthor={BlackArch Linux, BlackArch Linux},
pdfkeywords={BlackArch Linux, Penetration Testing, Security, Hacking, Linux},
pdfcenterwindow=true
]{hyperref}


%%% LAYOUT %%%
\setlength{\parindent}{0em}
\setlength{\parskip}{1.5ex plus0.5ex minus0.5ex}
\geometry{left=2.5cm, textwidth=16cm, top=3cm, textheight=25cm, bottom=3cm}
\widowpenalty=2000
\clubpenalty=1000
\frenchspacing
\sloppy
\pagecolor{black}
\color{white}
\hypersetup{pdftex=true, colorlinks=true, breaklinks=true, linkcolor=red,
menucolor=red, pagecolor=red, urlcolor=red}
\setcounter{tocdepth}{10}
\setcounter{secnumdepth}{10}


%%% HEADER / FOOTER %%%
\setlength{\headheight}{1.0cm}
\setlength{\headsep}{1.0cm}
\lhead{{\includegraphics[width=1cm,height=1cm]{images/logo.png}}}
\rhead{The BlackArch Linux Guide}
%\lfoot{\small{foo}}
%\rfoot{\small{bar}}


%%% OWN MACROS %%%
% put own macros here, like renewcommand newcommand etc.


%%% DOCUMENT %%%
\begin{document}
\pagestyle{empty}
\begin{center}
\begin{figure}[htbp]
\centering
\vspace{1cm}
\Huge{\textbf{BlackArch \color{red}Linux}}\\
\vspace{2cm}
\includegraphics[width=8cm]{images/logo.png}
\label{fig:logo}
\end{figure}
\vspace{1cm}
\Huge{\textbf{Guide de BlackArch}}\\
\vspace{1cm}
\Large{\color{red}http://www.blackarch.org/}\\
\vspace{0.5cm}
\Large{\today}
\end{center}
\newpage
\tableofcontents
\newpage
\pagestyle{fancy}

%%%%%%%%%%%%%%%%%%%%%%%%%%%%%%%%%%%%%%%%%%%%%%%%%%%%%%%%%%%%%%%%%%%%%%%%%%%%%%%%

\chapter{Introduction}

\section{Qu'est-ce que BlackArch Linux?}
\href{http://www.blackarch.org}{BlackArch Linux} est une légère extension d'Arch
Linux destinée aux professionnels de la sécurité informatique.
\\\\
L'ensemble d'outils est distribué tel un
\href{https://wiki.archlinux.org/index.php/Unofficial\_User\_Repositories}
{dépôt non-officiel} d'Arch Linux, vous permettant d'installer les composantes de
BlackArch sur une installation d'Arch Linux existante. Les paquets peuvent y être
installés individuellement ou par catégorie.
\\\\
Le dépôt logiciel de BlackArch contient plus de 1500 outils et ce nombre augmente
sans cesse. Tous ces outils sont minutieusement testés avant d'être ajoutés aux
dépôts, afin de garantir leur bon fonctionnement.

%\section{The story of BlackArch Linux}
%foo bar

%\section{Supported platforms}
%foo bar

\section{S'impliquer}
Il est possible d'entrer en contact avec l'équipe BlackArch à l'aide de ces
informations:
\\\\
Web: \url{http://www.blackarch.org/}
\\\\
Courriel: \href{mailto:blackarchlinux@gmail.com}{blackarchlinux@gmail.com}
\\\\
IRC: \url{irc://irc.freenode.net/blackarch}

%%%%%%%%%%%%%%%%%%%%%%%%%%%%%%%%%%%%%%%%%%%%%%%%%%%%%%%%%%%%%%%%%%%%%%%%%%%%%%%%

\chapter{Guide d'utilisation}

\section{Installation}
La section suivante montre comment configurer le dépôt de paquets de BlackArch ainsi
que la procédure à suivre afin d'installer des logiciels à partir de
celle-ci. BlackArch supporte l'installation de fichiers binaires précompilés
ainsi que l'installation à l'aide du code source original des paquets.
\\\\
BlackArch est compatible avec une installation de base d'Arch Linux. Il s'agira
d'un dépôt non-officiel de paquets. Si vous désirez une image ISO à la place,
visitez la section \href{http://www.blackarch.org/download.html#iso}{Live ISO}.
\\\\

\subsection{Installer le dépôt}
Il suffit de 6 étapes afin d'installer le dépôt de BlackArch sur une installation
existante de Arch Linux. Vous devez suivre les instructions dans l'ordre. N'ajoutez pas \textbf{blackarch}
à votre fichier \textit{pacman.conf} sans avoir complété les étapes 0 à 2.
\begin{enumerate} \itemsep4pt
\item Si vous avez déjà installé BlackArch et qu'il y existe déjà une entrée
\textbf{[blackarch]} dans \textit{/etc/pacman.conf}, vous devrez la supprimer ou
la mettre en commentaire, et ensuite exécuter la commande \textit{pacman -Syy}.
\item Exécutez les commandes suivantes en tant que root. Ces commande
téléchargeront et installeront les clés nécessaires pour garantir l'authenticité
des paquets via une signature digitale.
{\small
\color{gray}
\begin{verbatim}
wget -q http://blackarch.org/keyring/blackarch-keyring.pkg.tar.xz{,.sig}
gpg --keyserver hkp://pgp.mit.edu --recv 4345771566D76038C7FEB43863EC0ADBEA87E4E3
gpg --keyserver-o no-auto-key-retrieve --with-f blackarch-keyring.pkg.tar.xz.sig
pacman-key --init
rm blackarch-keyring.pkg.tar.xz.sig
pacman --noc -U blackarch-keyring.pkg.tar.xz
\end{verbatim}
}
\item (Facultatif) Vérifiez l'authenticité des clés d'autant de sources que vous
désirez, afin de vous convaincre de leur validité.
\item Ajoutez les lignes suivantes à votre fichier \textit{/etc/pacman.conf}:
{\small
\color{gray}
\begin{verbatim}
[blackarch]
Server = <site_miroir>/$repo/os/$arch
\end{verbatim}
}
Remplacez \textit{\textless site\_mirroir\textgreater} par le miroir officiel
de votre choix. La liste des miroirs est disponible sur
\href{http://www.blackarch.org/}{notre site web}.

\item Exécutez:
{\small
\color{gray}
\begin{verbatim}
# pacman -Syyu
\end{verbatim}
}
\end{enumerate}

\subsection{Installation de paquets}
Vous pouvez désormais installer des outils directement depuis le dépôt de BlackArch.
\begin{enumerate}
\item Afin d'obtenir la liste de tous les outils disponibles exécutez
{\small
\color{gray}
\begin{verbatim}
# pacman -Sgg | grep blackarch | cut -d' ' -f2 | sort -u
\end{verbatim}
}
\item Afin d'installer tous les outils, exécutez
{\small
\color{gray}
\begin{verbatim}
# pacman -S blackarch
\end{verbatim}
}
\item Afin d'installer une catégorie complète d'outils, exécutez
{\small
\color{gray}
\begin{verbatim}
# pacman -S blackarch-<nom de la catégorie>
\end{verbatim}
}
\item Afin d'obtenir la liste des catégories de paquets disponibles, exécutez
{\small
\color{gray}
\begin{verbatim}
# pacman -Sg | grep blackarch
\end{verbatim}
}
\end{enumerate}

\subsection{Installation des paquets depuis le code source}
Il existe une méthode alternative à l'installation de paquets sous forme binaire.
BlackArch vous offre la possibilité d'obtenir le code source directement depuis le dépôt officiel du paquet désiré, et BlackArch le compilera pour vous. Vous
trouverez les fichiers PKGBUILD sur notre
\href{https://github.com/BlackArch/blackarch/tree/master/packages}{github}. Afin
d'obtenir le dépôt en entier, vous pourrez utiliser l'outil
\href{https://github.com/BlackArch/blackman}{blackman}.
\begin{itemize} \itemsep10pt
\item Vous devez premièrement installer blackman. Si vous avez déjà configuré
le dépôt de paquets de BlackArch sur votre système, tel qu'indiqué au début de
ce document, vous pouvez installer blackman à l'aide de pacman:
{\small
\color{gray}
\begin{verbatim}
pacman -S blackman
\end{verbatim}
}
\item Vous pouvez aussi installer blackman directement depuis le code source:
{\small
\color{gray}
\begin{verbatim}
mkdir blackman
cd blackman
wget https://raw2.github.com/BlackArch/blackarch/master/packages/blackman/PKGBUILD
# Make sure the PKGBUILD has not been maliciously tampered with.
makepkg -s
\end{verbatim}
}
\item Afin de télécharger et compiler un paquet:
{\small
\color{gray}
\begin{verbatim}
$ sudo blackman -i package
\end{verbatim}
}
\item Afin de télécharger, compiler, et installer une catégorie entière:
{\small
\color{gray}
\begin{verbatim}
$ sudo blackman -g group
\end{verbatim}
}
\item Afin de télécharger, compiler, et installer tous les outils présents
dans le dépôt de paquets de BlackArch:
{\small
\color{gray}
\begin{verbatim}
$ sudo blackman -a
\end{verbatim}
}
\item Afin d'obtenir la liste de toutes les catégories de paquets de BlackArch:
{\small
\color{gray}
\begin{verbatim}
$ blackman -l
\end{verbatim}
}
\item Si vous voulez ensuite obtenir la liste de tous les paquets compris dans
une catégorie particulière, vous pouvez obtenir la liste de ceux-ci ainsi qu'une
courte description pour chaque paquet à l'aide de la commande suivante:
{\small
\color{gray}
\begin{verbatim}
$ blackman -p <nom de la catégorie>
\end{verbatim}
}
\end{itemize}

\subsection{Installation d'un LiveCD, netinstall, ou d'Arch Linux}
Vous pouvez installer BlackArch d'une de nos images ISO, que ce soit un LiveCD
ou un netcd. Vous les trouverez sur
\href{http://www.blackarch.org/download.html#iso}{notre site}. Les étapes
suivantes sont nécéssaires afin de créer une image bootable:
\begin{itemize}
\item Installez le paquet blackarch-installer:
{\small
\color{gray}
\begin{verbatim}
$ sudo pacman -S blackarch-installer
\end{verbatim}
}
\item Ensuite exécutez:
{\small
\color{gray}
\begin{verbatim}
$ sudo blackarch-install
\end{verbatim}
}
\end{itemize}

%%%%%%%%%%%%%%%%%%%%%%%%%%%%%%%%%%%%%%%%%%%%%%%%%%%%%%%%%%%%%%%%%%%%%%%%%%%%%%%%


\chapter{Guide de développement}

\section{Le Arch Build System et les dépôts de paquets}

Les fichiers PKGBUILD sont de simples scripts permettant d'installer et de
compiler un programme sur plusieurs environnements différents.
Chaque PKGBUILD indique à la commande makepkg(1) comment créer un paquet valide
qui pourra ensuite être installé. Les fichiers PKGBUILD sont écrits en bash.

Pour plus d'informations, il est suggéré de lire ces pages:
\begin{itemize}
\item \href{http://wiki.archlinux.fr/Standard_paquetage}{Arch Wiki FR:
Standard paquetage}
\item \href{http://wiki.archlinux.fr/Makepkg}{Arch Wiki FR: Makepkg}
\item \href{http://wiki.archlinux.fr/PKGBUILD}{Arch Wiki FR: PKGBUILD}
\end{itemize}

\section{Le standard de paquetage de BlackArch}

Afin de garder les choses simples, nos PKGBUILDs sont similaires à ceux du
Arch User Repository, à quelques différences près. Tous les paquets doivent
être membre du groupe \textit{blackarch}. Il est possible qu'un paquet soit
membre de plus d'un groupe.

\subsection{Groupes}

Afin de permettre aux utilisateurs d'installer certaines catégories de paquets
exclusivement, ceux-ci ont étés répertoriés en plusieurs groupes. Ces groupes
permettent aux utilisateurs d'exécuter la commande
\textit{pacman -S \textless nom\_du\_groupe \textgreater } afin d'installer tous
 les paquets compris dans ce groupe.

\subsubsection{blackarch}

En théorie, le groupe \textit{blackarch} doit contenir tous les paquets qu'offre
BlackArch. La majorité des paquets devrait être membre de ce groupe.

\subsubsection{blackarch-anti-forensic}

Groupe contenant les paquets utilisés pour contrer les tentatives
d'investigation informatique de systèmes. Vous y trouverez des paquets offrant
des solutions de cryptage de données, stéganographie, modification d'attributs
de fichiers, et tout autre type de logiciels permettant la modification d'un
système de fichiers dans le but d'y masquer de l'information.

Exemples: luks, TrueCrypt, Timestomp, dd, ropeadope, secure-delete

\subsubsection{blackarch-automation}

Groupe contenant les paquets utilisés pour l'automatisation de plusieurs tâches.
Puisque les paquets présents dans ce groupe varient énormément en
fonctionnalité, il vous est suggéré de consulter directement les paquets
afin de vous informer sur leurs capacités.

Exemples: blueranger, tiger, wiffy

\subsubsection{blackarch-backdoor}

Groupe contenant les paquets liés à l'exploitation ou à l'ouverture de
backdoors sur des systèmes vulnérables.

Exemples: backdoor-factory, rrs, weevely

\subsubsection{blackarch-binary}

Groupe contenant les paquets reliés aux fichiers binaires. Vous y trouverez des
paquets permettant l'extraction d'archives d'un firmware, des logiciels de
rétro-ingénierie, logiciels aidant à l'exploitation de fichiers binaires, etc.

Exemples: binwally, packerid, hex2bin, binwalk

\subsubsection{blackarch-bluetooth}

Groupe contenant les paquets aidant à l'exploitation du standard Bluetooth
(802.15.1).

Exemples: ubertooth, tbear, redfang

\subsubsection{blackarch-code-audit}

Groupe contenant les paquets à fin d'audit de code source. Ces paquets analysent
statiquement le code source de votre projet afin d'y déceler des vulnérabilités.

Exemples: flawfinder, pscan

\subsubsection{blackarch-cracker}

Groupe contenant les paquets utilisés afin de tester la robustesse de multiples
fonctions cryptographiques.

Exemples: hashcat, john, crunch

\subsubsection{blackarch-crypto}

Groupe contenant tous types de paquets ayant rapport de près ou de loin à la
cryptographie, que ce soit de l'analyse ou du craquage de hash.

Exemples: ciphertest, xortool, sbd

\subsubsection{blackarch-database}

Groupe contenant les paquets reliés à l'exploitation de bases de données.

Exemples: metacoretex, blindsql

\subsubsection{blackarch-debugger}

Groupe contenant les paquets permettant à l'usager d’investiguer le
fonctionnement interne d'un programme lors de son utilisation.

Exemples: radare2, shellnoob

\subsubsection{blackarch-decompiler}

Groupe contenant les paquets qui tentent d'effectuer de la rétro-ingénierie sur
un programme compilé.

Exemples: flasm, jd-gui

\subsubsection{blackarch-defensive}

Groupe contenant les paquets utilisés afin de protéger un système des logiciels
et usagers malveillants.

Exemples: arpon, chkrootkit, sniffjoke

\subsubsection{blackarch-disassembler}

Ce groupe est fortement similaire à \textit{blackarch-decompiler} dans le sens
où ces deux groupes tentent d'effectuer une rétro-ingénierie générique de
fichiers exécutables, mais les paquets compris dans
\textit{blackarch-disassembler} en extrairont le code assembleur au lieu d'un
pseudocode.

Exemples: inguma, radare2

\subsubsection{blackarch-dos}

Groupe contenant les paquets aidant aux attaques de type \textit{déni de service}.

Exemples: 42zip, nkiller2

\subsubsection{blackarch-drone}

Groupe contenant les paquets aidant à l'entretien et l'utilisation de drones.

Exemples: meshdeck, skyjack

\subsubsection{blackarch-exploitation}

Groupe contenant les paquets utilisés lors de l'exploitation de multiples
programmes et services.

Exemples: armitage, metasploit, zarp

\subsubsection{blackarch-fingerprint}

Groupe contenant les paquets utilisés à fin de reconnaissance réseautique.

Exemples: dns-map, p0f, httprint

\subsubsection{blackarch-firmware}

Groupe contenant les paquets utilisés lors de l'exploitation de firmware.

Exemples: None yet, amend asap.

\subsubsection{blackarch-forensic}

Groupe contenant les paquets utilisés lors d'investigation informatique portant
sur le recouvrement de données potentiellement effacés.

Exemples: aesfix, nfex, wyd

\subsubsection{blackarch-fuzzer}

Groupe contenant les paquets aidant au fuzzing.

Exemples: msf, mdk3, wfuzz

\subsubsection{blackarch-hardware}

Groupe contenant les paquets aidant à l'exploitation de matériel physique.

Exemples: arduino, smali

\subsubsection{blackarch-honeypot}

Groupe contenant les paquets agissant comme des \textit{honeypots}, c'est-à-dire
des services sécuritaires apparaissant vulnérables, afin d'attirer les \textit{hackers} dans un piège.

Exemples: artillery, bluepot, wifi-honey

\subsubsection{blackarch-keylogger}

Groupe contenant les paquets utilisés lors d'installation de \textit{keylogging}, 
soit l'enregistrement de touches clavier sur un système, afin d'en soutirer des informations confidentielles.

Exemples: klogger, logkeys, xspy

\subsubsection{blackarch-malware}

Groupe contenant tous types de paquets ayant un lien au \textit{malware}, que ce
soit l'utilisation de ceux-ci ou leur détection.

Exemples: malwaredetect, peepdf, yara

\subsubsection{blackarch-misc}

Groupe contenant les paquets semblant prendre part dans aucune autre
catégorie.

Exemples: oh-my-zsh-git, winexe, stompy

\subsubsection{blackarch-mobile}

Groupe contenant les paquets aidant aux activités d'entretien et d'audit
d'application, de développement d'application et de matériel mobile.

Exemples: android-sdk-platform-tools, android-udev-rules

\subsubsection{blackarch-networking}

Groupe contenant les paquets aidant à de multiples tâches liées à la réseautique.

Exemples: Anything pretty much

\subsubsection{blackarch-nfc}

Groupe contenant les paquets utilisant la technologie \textit{NFC} (Near-Field
Communications).

Exemples: nfcutils

\subsubsection{blackarch-packer}

Groupe contenant les paquets utilisant de près ou de loin les \textit{packers},
soit des logiciels qui créent des applications contenant un programme
malveillant caché à l'intérieur.

Exemples: packerid

\subsubsection{blackarch-proxy}

Groupe contenant les paquets agissant tel un proxy, soit un logiciel redirigeant
le trafic internet afin de l'analyser ou de le modifier.

Exemples: burpsuite, ratproxy, sslnuke

\subsubsection{blackarch-recon}

Groupe contenant les paquets servant à la recherche de services exploitables dans
le monde réel.

Exemples: canri, dnsrecon, netmask

\subsubsection{blackarch-reversing}

Groupe contenant les paquets de décompilation, de désassemblage, et tout autre
type de programme ayant à voir avec la rétro-ingénierie.

Exemples: capstone, radare2, zerowine

\subsubsection{blackarch-scanner}

Groupe contenant les paquets utilisés afin de scanner de multiples systèmes pour 
y déceler des vulnérabilités actives.

Exemples: scanssh, tiger, zmap

\subsubsection{blackarch-sniffer}

Groupes contenant les paquets aidant à l'analyse de trafic réseau.

Exemples: hexinject, pytactle, xspy

\subsubsection{blackarch-social}

Groupe contenant les paquets attaquant principalement les sites de réseaux
sociaux.

Exemples: jigsaw, websploit

\subsubsection{blackarch-spoof}

Groupe contenant les paquets qui tentent de dissimuler un attaquant, pour que 
la victime ne puisse le percevoir comme un acteur malicieux.

Exemples: arpoison, lans, netcommander

\subsubsection{blackarch-threat-model}

Groupe contenant les paquets utilisés pour enregistrer et créer des rapports sur
plusieurs modèles de menaces pouvant se produire lors d'un scénario spécifique.

Exemples: magictree

\subsubsection{blackarch-tunnel}

Groupe contenant les paquets utilisés pour encapsuler le contenu réseau via un
tiers parti, afin de dissimuler la provenance initiale de données.

Exemples: ctunnel, iodine, ptunnel

\subsubsection{blackarch-unpacker}

Groupe contenant les paquets utilisés lors de l'extraction de logiciel
malveillants à partir d'un éxécutable.

Exemples: js-beautify

\subsubsection{blackarch-voip}

Groupe contenant les paquets utilisés afin d'effectuer des actions ou de l'analyse
de téléphonie VoIP.

Exemples: iaxflood, rtp-flood, teardown

\subsubsection{blackarch-webapp}

Groupe contenant les paquets effectuant des opérations de validation de sécurité
sur des applications accessibles via l'internet, afin d'en garantir leur
robustesse.

Exemples: metoscan, whatweb, zaproxy

\subsubsection{blackarch-windows}

Groupe contenant les paquets Windows qui s'exécuteront via Wine.

Exemples: 3proxy-win32, pwdump, winexe

\subsubsection{blackarch-wireless}

Groupe contenant les paquets utilisés lors d'audits de sécurité sur une
infrastructure sans-fil.

Exemples: airpwn, mdk3, wiffy

\section{Structure du dépôt}

Vous pouvez trouver le dépôt principal de BlackArch sur
\href{https://github.com/BlackArch/blackarch}
{https://github.com/BlackArch/blackarch}. Il existe aussi plusieurs dépôts de
code sur \href{https://github.com/BlackArch}{https://github.com/BlackArch}.

À l'intérieur du dépôt principal, vous trouverez les trois dossiers principaux:

\begin{itemize}
\item docs - Contient la documentation du projet
\item packages - Contient les fichiers PKGBUILD
\item scripts - Contient de multiples scripts d'automation
\end{itemize}

\subsection{Scripts}

Voici une référence des scripts que vous trouverez dans le dossier
\verb|scripts/| :

\begin{itemize}
\item baaur - Éventuellement, ce script enverra nos paquets à la banque de données
du AUR.
\item babuild - Construit un paquet.
\item bachroot - Prépare un chroot afin d'effectuer des tests.
\item baclean - Nettoie les anciens fichiers .pkg.tar.xz du dépôt.
\item baconflict - Éventuellement, ce script remplacera \verb|scripts/conflicts|.
\item bad-files - Trouve des fichiers erronés dans les paquets construits.
\item balock - Crée ou termine le lock sur le dépôt de paquets.
\item banotify - Indique IRC des push de paquets.
\item barelease - Envoie les paquets au dépôt.
\item baright - Affiche les informations de copyright de BlackArch.
\item basign - Signe les paquets.
\item basign-key - Signe une clé.
\item blackman - Opère sensiblement comme pacman mais construit les paquets à
partir de git. (à ne pas confondre avec le programme blackman, écrit par nzr)
\item check-groups - Valide les groupes.
\item checkpkgs - Valide les paquets.
\item conflicts - Valide les paquets pour des conflits.
\item dbmod - Modifie la base de données d'un paquet.
\item depth-list - Crée une liste de dépendances ordonnée par profondeur.
\item deptree - Crée un arbre de dépendances, ne listant que les paquets obtenus
via BlackArch.
\item get-blackarch-deps - Affiche une liste des dépendances de paquets obtenues
via BlackArch pour un paquet donné.
\item get-official - Obtient les paquets officiels.
\item list-loose-packages - Liste les paquets qui ne se trouvent dans aucun
groupe et ne sont pas des dépendances pour d'autres paquets.
\item list-needed - Liste les dépendances manquantes.
\item list-removed - Liste les paquets qui sont dans le dépôt mais pas dans le
répertoire git.
\item list-tools - Liste les outils.
\item outdated - Cherche les paquets qui ne sont plus à jour lorsque comparés
avec le répertoire git.
\item pkgmod - Modifie un paquet de construction.
\item pkgrel - Incrémente le pkgrel d'un paquet.
\item prep - Nettoie un fichier PKGBUILD de ses erreurs de style et autre.
\item sitesync - Synchronise entre une copie locale et sa version distante.
\item size-hunt - Cherche de gros paquets.
\item source-backup - Sauvegarde les fichiers source d'un paquet.
\end{itemize}

\section{Contribuer au dépôt}
Cette section indique comment contribuer au projet BlackArch Linux. Nous
acceptons les pull requests de toute taille, que ce soit de petites
modifications ou un ajout de paquet.\\Pour recevoir de l'aide, n'hésitez pas à
nous contacter.
\\\\
Tout le monde est le bienvenu pour contribuer. Toute contribution sera appréciée.

\subsection{Tutoriels requis}
Veuillez lire les documents suivants avant de contribuer:
\begin{itemize}
\item \href{http://wiki.archlinux.fr/Standard_paquetage}{Arch Wiki FR:
Standard paquetage}
\item \href{http://wiki.archlinux.fr/Makepkg}{Arch Wiki FR: Makepkg}
\item \href{http://wiki.archlinux.fr/PKGBUILD}{Arch Wiki FR: PKGBUILD}
\end{itemize}

\subsection{Étapes à suivre afin de contribuer}
Afin de contribuer au projet BlackArch Linux, suivez ces étapes:
\begin{enumerate}
\item Forkez le dépôt sur github à partir de
\url{https://github.com/BlackArchLinux/blackarchlinux}.
\item Effectuez vos modifications sur les fichiers nécéssaires (e.g. PKGBUILD,
.patch files, etc).
\item Commitez vos changements.
\item Poussez vos changements.
\item Créez un pull request vers le dépôt.
\end{enumerate}

\subsection{Exemple}
L'exemple suivant montre la soumission d'un nouveau paquet au projet
BlackArch. Nous utilisons \href{https://wiki.archlinux.org/index.php/yaourt}{yaourt}
(vous pouvez aussi utiliser pacaur) afin d'obtenir le PKGBUILD existant de
\textbf{nfsshell}, et nous pouvons ensuite l'éditer, au besoin.

\subsubsection{Obtenir le PKGBUILD}
En utilisant yaourt ou pacaur, obtenir \textit{PKGBUILD}:
{\small
\color{gray}
\begin{verbatim}
user@blackarchlinux $ yaourt -G nfsshell
==> Download nfsshell sources
x LICENSE
x PKGBUILD
x gcc.patch
user@blackarchlinux $ cd nfsshell/
\end{verbatim}
}

\subsubsection{Nettoiement du PKGBUILD}
Nettoyer le fichier \textit{PKGBUILD} à l'aide de \textit{prep}, afin de sauver
du temps:
{\small
\color{gray}
\begin{verbatim}
user@blackarchlinux nfsshell $ ./blarckarch/scripts/prep PKGBUILD
cleaning 'PKGBUILD'...
expanding tabs...
removing vim modeline...
removing id comment...
removing contributor and maintainer comments...
squeezing extra blank lines...
removing '|| return'...
removing leading blank line...
removing $pkgname...
removing trailing whitespace...
\end{verbatim}
}

\subsubsection{Ajustement du PKGBUILD}
Ajustez le fichier \textit{PKGBUILD}:
{\small
\color{gray}
\begin{verbatim}
user@blackarchlinux nfsshell $ vi PKGBUILD
\end{verbatim}
}

\subsubsection{Construire le paquet}
Construire le paquet:
{\small
\color{gray}
\begin{verbatim}
user@blackarchlinux nfsshell $ makepkg -sf
==> Making package: nfsshell 19980519-1 (Mon Dec 2 17:23:51 CET 2013)
==> Checking runtime dependencies...
==> Checking buildtime dependencies...
==> Retrieving sources...
-> Downloading nfsshell.tar.gz...
% Total % Received % Xferd Average Speed Time Time Time
CurrentDload Upload Total Spent Left Speed100 29213 100 29213 0
0 48150 0 --:--:-- --:--:-- --:--:-- 48206
-> Found gcc.patch
-> Found LICENSE
...
<beaucoup d'output du processus de build et compilation ici>
...
==> Leaving fakeroot environment.
==> Finished making: nfsshell 19980519-1 (Mon Dec 2 17:23:53 CET 2013)
\end{verbatim}
}

\subsubsection{Installation et test du paquet}
Installation et test du paquet:
{\small
\color{gray}
\begin{verbatim}
user@blackarchlinux nfsshell $ pacman -U nfsshell-19980519-1-x86_64.pkg.tar.xz
user@blackarchlinux nfsshell $ nfsshell # test it
\end{verbatim}
}

\subsubsection{Ajout, commit et push du paquet}
Ajout, commit et push du paquet:
{\small
\color{gray}
\begin{verbatim}
user@blackarchlinux nfsshell $ cd /blackarchlinux/packages
user@blackarchlinux ~/blackarchlinux/packages $ mv ~/nfsshell .
user@blackarchlinux ~/blackarchlinux/packages $ git commit -am nfsshell && git push
\end{verbatim}
}

\subsubsection{Création d'un pull request}
Créez un pull request de votre fork vers
 \href{https://github.com/BlackArchLinux/blackarchlinux}{le projet officiel}:
{\small
\color{gray}
\begin{verbatim}
firefox https://github.com/<contributor>/blackarchlinux
\end{verbatim}
}

\subsubsection{Ajout d'un remote pour l'upstream}
Une fois votre modification faite, il serait une bonne idée d'ajouter le projet
officiel comme remote, afin d'ajouter les modifications du projet à votre fork.
Cela vous permettra de synchroniser votre fork avec BlackArch.
{\small
\color{gray}
\begin{verbatim}
user@blackarchlinux ~/blackarchlinux $ git remote -v
origin <the url of your fork> (fetch)
origin <the url of your fork> (push)
user@blackarchlinux ~/blackarchlinux $ git remote add upstream https://github.com/blackarch/blackarch
user@blackarchlinux ~/blackarchlinux $ git remote -v
origin <the url of your fork> (fetch)
origin <the url of your fork> (push)
upstream https://github.com/blackarch/blackarch (fetch)
upstream https://github.com/blackarch/blackarch (push)
\end{verbatim}
}
Par défaut, git devrait push directement à l'origine, mais assurez-vous que
votre configuration git soit correcte. Cela ne sera pas un problème à moins
d'avoir les droits de commit puisque vous ne pourrez pas push upstream sans
cela.

Si vous avez les droits de commit, vous aurez possiblement plus de succès en
utilisant \textit{git@github.com:blackarch/blackarch.git}. Le choix est vôtre.

\subsection{Requêtes}
\begin{enumerate}
\item N'ajoutez pas de commentaire \textbf{Maintainer} ou \textbf{Contributor}
au fichiers \textit{PKGBUILD}. Ajoutez ces commentaires à la section
\textbf{AUTHORS} du guide BlackArch.
\item Afin de conserver la consistance du projet, veuillez suivre les règles
générales de style des autres fichiers \textit{PKGBUILD} en utilisant une
indentation de deux espaces.
\end{enumerate}

\subsection{Conseils généraux}
\href{http://wiki.archlinux.org/index.php/Namcap}{namcap} peut trouver des
erreurs dans vos paquets. Il est recommandé de l'utiliser avant de soumettre vos
modifications.

%%% APPENDIX %%%
\appendix
%%%%%%%%%%%%%%%%%%%%%%%%%%%%%%%%%%%%%%%%%%%%%%%%%%%%%%%%%%%%%%%%%%%%%%%%%%%%%%%%
%                                                                              %
% BlackArch Linux plantilla de ap\'endice                                      %
%                                                                              %
%%%%%%%%%%%%%%%%%%%%%%%%%%%%%%%%%%%%%%%%%%%%%%%%%%%%%%%%%%%%%%%%%%%%%%%%%%%%%%%%

\appendix

\chapter{Ap\'endice}

\section{FAQ's}

\section{AUTORES}
Las siguientes personas han colaborado directamente en BlackArch:
\begin{itemize}
\item Tyler Bennnett (tylerb@trix2voip.com)
\item fnord0 (fnord0@riseup.net)
\item nrz (nrz@nullsecurity.net)
\item Ellis Kenyo (elken.tdos@gmail.com)
\item CaledoniaProject (the.warl0ck.1989@gmail.com)
\item sudokode (sudokode@gmail.com)
\item Valentin Churavy (v.churavy@gmail.com)
\item Boy Sandy Gladies Arriezona (reno.esper@gmail.com)
\item Mathias Nyman
\item Johannes Löthberg (demizide@gmail.com)
\item Thiago da Silva Teixeira (teixeira.zeus@gmail.com)
\end{itemize}
Las siguientes personas han colaborado directamente en ArchPwn, 
y se han unido a BlackArch:
\begin{itemize}
\item Francesco Piccinno (stack.box@gmail.com)
\item jensp (jens@jenux.homelinux.org)
\item Valentin Churavy (v.churavy@gmail.com)
\end{itemize}
Hemos tomado build code de las siguientes personas:
\begin{itemize}
\item  3ED (krzysztof1987@gmail.com)
\item  AUR Perl (aurperl@juster.info)
\item  Aaron Griffin (aaron@archlinux.org)
\item  Abakus (java5@arcor.de)
\item  Adam Wolk (netprobe@gmail.com)
\item  Aleix Pol (aleixpol@kde.org)
\item  Aleshus (aleshusi@gmail.com)
\item  Alessandro Pazzaglia (jackdroido@gmail.com)
\item  Alessandro Sagratini (ale\_sagra@hotmail.com)
\item  Alex Cartwright (alexc223@googlemail.com)
\item  Alexander De Sousa (archaur.xandy21@spamgourmet.com)
\item  Alexander Rødseth (rodseth@gmail.com)
\item  Allan McRae (allan@archlinux.org)
\item  AmaN (gabroo.punjab.da@gmail.com)
\item  Andre Klitzing (aklitzing@online.de)
\item  Andrea Scarpino (andrea@archlinux.org)
\item  Andreas Schönfelder (passtschu@freenet.de)
\item  Andrej Gelenberg (andrej.gelenberg@udo.edu)
\item  Angel Velasquez (angvp@archlinux.org)
\item  Antoine Lubineau (antoine@lubignon.info)
\item  Anton Bazhenov (anton.bazhenov@gmail.com)
\item  Arkham (arkham@archlinux.us)
\item  Arthur Danskin (arthurdanskin@gmail.com)
\item  Balda (balda@balda.ch)
\item  Balló György (ballogyor+arch@gmail.com)
\item  Bartek Piotrowski (barthalion@gmail.com)
\item  Bartosz Feński (fenio@debian.org)
\item  Bartłomiej Piotrowski (nospam@bpiotrowski.pl)
\item  Bogdan Szczurek (thebodzio@gmail.com)
\item  Brad Fanella (bradfanella@archlinux.us)
\item  Brian Bidulock (bidulock@openss7.org)
\item  C Anthony Risinger (anthony@xtfx.me)
\item  CRT (crt.011@gmail.com)
\item  Can Celasun (dcelasun@gmail.com)
\item  Chaniyth (chaniyth@yahoo.com)
\item  Chris Brannon (cmbrannon79@gmail.com)
\item  Chris Giles (Chris.G.27@gmail.com) \& daschu117
\item  Christoph Siegenthaler (csi@gmx.ch)
\item  Christoph Zeiler (archNOSPAM@moonblade.org)
\item  Clément DEMOULINS (clement@archivel.fr)
\item  Corrado Primier (bardo@aur.archlinux.org)
\item  Daenyth (Daenyth+Arch@gmail.com)
\item  Dale Blount (dale@archlinux.org)
\item  Damir Perisa (damir.perisa@bluewin.ch)
\item  Dan Fuhry (dan@fuhry.us)
\item  Dan Serban (dserban01@yahoo.com)
\item  Daniel A. Campoverde Carrión
\item  Daniel Golle
\item  Daniel Griffiths (ghost1227@archlinux.us)
\item  Daniel J Griffiths (ghost1227@archlinux.us)
\item  Daniel Micay (danielmicay@gmail.com)
\item  Dave Reisner (dreisner@archlinux.org)
\item  Dawid Wrobel (cromo@klej.net)
\item  Devaev Maxim (mdevaev@gmail.com)
\item  Devin Cofer (ranguvar@archlinux.us)
\item  DigitalPathogen (aur@InfoSecResearchLabs.co.uk)
\item  DigitalPathogen (aur@digitalpathogen.co.uk)
\item  Dmitry A. Ilyashevich (dmitry.ilyashevich@gmail.com)
\item  Dominik Heidler (dheidler@gmail.com)
\item  DrZaius (lou@fakeoutdoorsman.com)
\item  Ebubekir KARUL (ebubekirkarul@yandex.com)
\item  Eduard "bekks" Warkentin (eduard.warkentin@gmail.com)
\item  Elmo Todurov (todurov@gmail.com)
\item  Emmanuel Gil Peyrot (linkmauve@linkmauve.fr)
\item  Eric Belanger (eric@archlinux.org)
\item  Ermak (ermak@email.it)
\item  Evangelos Foutras (evangelos@foutrelis.com)
\item  Fabian Melters (melters@gmail.com)
\item  Fabiano Furtado (fusca14@gmail.com)
\item  Federico Quagliata (ntp@quaqo.org)
\item  Firmicus (francois.archlinux@org)
\item  Florian Pritz (bluewind@jabber.ccc.de)
\item  Florian Pritz (flo@xinu.at)
\item  Francesco Piccinno (stack.box@gmail.com)
\item  François Charette (francois@archlinux.org)
\item  Gaetan Bisson (bisson@archlinux.org)
\item  Geoffroy Carrier (geoffroy.carrier@koon.fr)
\item  Georg Grabler (STiAT)
\item  George Hilliard (gh403@msstate.edu)
\item  Gerardo Exequiel Pozzi (vmlinuz386@yahoo.com.ar)
\item  Gilles CHAUVIN (gcnweb@gmail.com)
\item  Giovanni Scafora (giovanni@archlinux.org)
\item  Gordin (9ordin@gmail.com)
\item  Guillaume ALAUX (guillaume@archlinux.org)
\item  Guillermo Ramos (0xwille@gmail.com)
\item  Gustavo Alvarez (sl1pkn07@gmail.com)
\item  Hugo Doria (hugo@archlinux.org)
\item  Hyacinthe Cartiaux (hyacinthe.cartiaux@free.fr)
\item  James Fryman (jfryman@gmail.com)
\item  Jan "heftig" Steffens (jan.steffens@gmail.com)
\item  Jan de Groot (jgc@archlinux.org)
\item  Jaroslav Lichtblau (dragonlord@aur.archlinux.org)
\item  Jaroslaw Swierczynski (swiergot@aur.archlinux.org)
\item  Jason Chu (jason@archlinux.org)
\item  Jason R Begley (jayray@digitalgoat.com)
\item  Jason Rodriguez
\item  Jason St. John (jstjohn@purdue.edu)
\item  Jawmare (victor2008@gmail.com)
\item  Jeff Mickey (jeff@archlinux.org)
\item  Jens Pranaitis (jens@chaox.net)
\item  Jens Pranaitis (jens@jenux.homelinux.org)
\item  Jinx (jinxware@gmail.com)
\item  John D Jones III (jnbek1972@gmail.com)
\item  John Proctor (jproctor@prium.net)
\item  Jon Bergli Heier (snakebite@jvnv.net)
\item  Jonas Heinrich
\item  Jonathan Steel (jsteel@aur.archlinux.org)
\item  Joris Steyn (jorissteyn@gmail.com)
\item  Josh VanderLinden (arch@cloudlery.com)
\item  Jozef Riha (jose1711@gmail.com)
\item  Judd Vinet (jvinet@zeroflux.org)
\item  Juergen Hoetzel (jason@archlinux.org)
\item  Juergen Hoetzel (juergen@archlinux.org)
\item  Justin Davis (jrcd83@gmail.com)
\item  Kaiting Chen (kaitocracy@gmail.com)
\item  Kaos
\item  Kevin Piche (kevin@archlinux.org)
\item  Kory Woods (kory@virlo.net)
\item  Kyle Keen (keenerd@gmail.com)
\item  Larry Hajali (larryhaja@gmail.com)
\item  LeCrayonVert
\item  Le\_suisse (lesuisse.dev+aur@gmail.com)
\item  Lekensteyn (lekensteyn@gmail.com)
\item  Limao Luo (luolimao+AUR@gmail.com)
\item  Lucien Immink
\item  Lukas Fleischer (archlinux@cryptocrack.de)
\item  Manolis Tzanidakis
\item  Marcin "avalan" Falkiewicz (avalatron@gmail.com)
\item  Mariano Verdu (verdumariano@gmail.com)
\item  Marti Raudsepp (marti@juffo.org)
\item  MatToufoutu (mattoufootu@gmail.com)
\item  Matthew Sharpe (matt.sharpe@gmail.com)
\item  Mauro Andreolini (mauro.andreolini@unimore.it)
\item  Max Pray a.k.a. Synthead (synthead@gmail.com)
\item  Max Roder (maxroder@web.de)
\item  Maxwell Pray a.k.a. Synthead (synthead@gmail.com)
\item  Maxwell Pray a.k.a. Synthead (synthead1@gmail.com)
\item  Mech (tiago.bmp@gmail.com)
\item  Michael Düll (mail@akurei.me)
\item  Michael P (ptchinster@archlinux.us)
\item  Michal Krenek (mikos@sg1.cz)
\item  Michal Zalewski (lcamtuf@coredump.cx)
\item  Miguel Paolino (mpaolino@gmail.com)
\item  Miguel Revilla (yo@miguelrevilla.com)
\item  Mike Roberts (noodlesgc@gmail.com)
\item  Mike Sampson (mike@sambodata.com)
\item  Nassim Kacha (nassim.kacha@gmail.com)
\item  Nicolas Pouillard (nicolas.pouillard@gmail.com)
\item  Nicolas Pouillard https://nicolaspouillard.fr
\item  Niklas Schmuecker
\item  Oleander Reis (oleander@oleander.cc)
\item  Olivier Le Moal (mail@olivierlemoal.fr)
\item  Olivier Médoc "oliv" (o\_medoc@yahoo.fr)
\item  Pascal E. (archlinux@hardfalcon.net)
\item  Patrick Leslie Polzer (leslie.polzer@gmx.net)
\item  Paul Mattal (paul@archlinux.org)
\item  Paul Mattal (pjmattal@elys.com)
\item  Pengyu CHEN (cpy.prefers.you@gmail.com)
\item  Peter Wu (lekensteyn@gmail.com)
\item  Philipp 'TamCore' B. (philipp@tamcore.eu)
\item  Pierre Schmitz (pierre@archlinux.de)
\item  Pranay Kanwar (pranay.kanwar@gmail.com)
\item  Pranay Kanwar (warl0ck@metaeye.org)
\item  PyroPeter (googlemail@com.abi1789)
\item  PyroPeter (googlemail.com@abi1789)
\item  Ray Rashif (schiv@archlinux.org)
\item  Remi Gacogne
\item  Renan Fernandes (renan@kauamanga.com)
\item  Richard Murri (admin@richardmurri.com)
\item  Roberto Alsina (ralsina@kde.org)
\item  Robson Peixoto (robsonpeixoto@gmail.com)
\item  Roel Blaauwgeers (roel@ttys0.nl)
\item  Rorschach (r0rschach@lavabit.com)
\item  Ruben Schuller (shiml@orgizm.net)
\item  Rudy Matela (rudy@matela.com)
\item  Ryon Sherman (ryon.sherman@gmail.com)
\item  Sabart Otto \item  Seberm (seberm@gmail.com)
\item  SakalisC (chrissakalis@gmail.com)
\item  Sam Stuewe (halosghost@archlinux.info)
\item  SanskritFritz (SanskritFritz@gmail.com)
\item  Sarah Hay (sarahhay@mb.sympatico)
\item  Sebastian Benvenuti (sebastianbenvenuti@gmail.com)
\item  Sebastian Nowicki (sebnow@gmail.com)
\item  Sebastien Duquette (ekse.0x@gmail.com)
\item  Sebastien LEDUC (sebastien@sleduc.fr)
\item  Sebastien Leduc (sebastien@sleduc.fr)
\item  Sergej Pupykin (pupykin.s+arch@gmail.com)
\item  Sergio Rubio (rubiojr@biondofu.net)
\item  Sheng Yu (magicfish1990@gmail.com)
\item  Simon Busch (morphis@gravedo.de)
\item  Simon Legner (Simon.Legner@gmail.com)
\item  Sirat18 (aur@sirat18.de)
\item  SpepS (dreamspepser@yahoo.it)
\item  Spider.007 (archPackage@spider007.net)
\item  Stefan Seering
\item  Stephane Travostino (stephane.travostino@gmail.com)
\item  Stéphane Gaudreault (stephane@archlinux.org)
\item  Sven Kauber (celeon@gmail.com)
\item  Sven Schulz (omee@archlinux.de)
\item  Sébastien Duquette (ekse.0x@gmail.com)
\item  Sébastien Luttringer (seblu@archlinux.org)
\item  TDY (tdy@gmx.com)
\item  Teemu Rytilahti (tpr@iki.fi)
\item  Testuser\_01
\item  Thanx (thanxm@gmail.com)
\item  Thayer Williams (thayer@archlinux.org)
\item  Thomas S Hatch (thatch45@gmail.com)
\item  Thorsten Töpper
\item  Tilmann Becker (tilmann.becker@freenet.de)
\item  Timothy Redaelli (timothy.redaelli@gmail.com)
\item  Timothée Ravier (tim@siosm.fr)
\item  Tino Reichardt
\item  Tobias Kieslich (tobias@justdreams.de)
\item  Tobias Powalowski (tpowa@archlinux.org)
\item  Tom K (tomk@runbox.com)
\item  Tom Newsom (Jeepster@gmx.co.uk)
\item  Tomas Lindquist Olsen (tomas@famolsen.dk)
\item  Travis Willard (travisw@wmpub.ca)
\item  Valentin Churavy (v.churavy@gmail.com)
\item  ViNS (gladiator@fastwebnet.it)
\item  Vlatko Kosturjak (kost@linux.hr)
\item  Wes Brown (wesbrown18@gmail.com)
\item  William Rea (sillywilly@gmail.com)
\item  Xavier Devlamynck (magicrhesus@ouranos.be)
\item  Xiao\item Long Chen (chenxiaolong@cxl.epac.to)
\item  aeolist (aeolist@hotmail.com)
\item  ality@pbrane.org
\item  astaroth (astaroth\_@web.de)
\item  bender02@archlinux.us
\item  billycongo (billycongo@gmail.com)
\item  bslackr (brendan@vastactive.com)
\item  cbreaker (cbreaker@tlen.pl)
\item  chimeracoder (dev@chimeracoder.net)
\item  damir (damir@archlinux.org)
\item  danitool
\item  darkapex (me@jailuthra.in)
\item  daronin
\item  dkaylor (dpkaylor@gmail.com)
\item  dobo (dobo90\_at\_gmail@com)
\item  dorphell (dorphell@archlinux.org)
\item  evr (evanroman@at.gmail)
\item  fnord0 (fnord0@riseup.net)
\item  fxbru (frxbru@gmail)
\item  hcar
\item  icarus (icarus.roaming@gmail.com)
\item  iceman (icemanf@gmail.com)
\item  kastor (kastor@fobos.org)
\item  kfgz (kfgz@interia.pl)
\item  linuxSEAT (linuxSEAT@gmail.com)
\item  m4xm4n (max@maxfierke.com)
\item  mar77i (mysatyre@gmail.com)
\item  marc0s (marc0s@fsfe.org)
\item  mickael9 (mickael9@gmail.com)
\item  nblock (nblock@archlinux.us)
\item  nofxx (x@nofxx.com)
\item  onny (onny@project
\item  pootzko (pootzko@gmail.com)
\item  revel (revel@muub.net)
\item  rich\_o (rich\_o@lavabit.com)
\item  s1gma (s1gma@mindslicer.com)
\item  sandman (r.coded@gmail.com)
\item  sebikul (sebikul@gmail.com)
\item  sh0 (mee@sh0.org)
\item  shild (sxp@bk.ru)
\item  simo (simo@archlinux.org)
\item  snuo
\item  sudokode (sudokode@gmail.com)
\item  tobias (tobias@archlinux.org)
\item  trashstar (trash@ps3zone.org)
\item  unexist (unexist@subforge.org)
\item  untitled (rnd0x00@gmail.com)
\item  virtuemood (virtue@archlinux.us)
\item  wido (widomaker2k7@gmail.com)
\item  wodim (neikokz@gmail.com)
\item  yannsen (ynnsen@gmail.com)
\end{itemize}


\end{document}

%%% EOF %%%
